%% ============================================================
%% manuscript.tex
%% Clinical Surgery Case-Time Prediction
%% Full LaTeX manuscript — compile with pdflatex (twice)
%% Figures generated by make_figures.py, flowchart by make_flowchart.py
%% ============================================================
\documentclass[11pt,twocolumn]{article}

%% ── Packages ─────────────────────────────────────────────────────────────────
\usepackage[T1]{fontenc}
\usepackage[utf8]{inputenc}
\usepackage{lmodern}
\usepackage[margin=1in,columnsep=0.35in]{geometry}
\usepackage{amsmath,amssymb}
\usepackage{graphicx}
\usepackage{booktabs}
\usepackage{multirow}
\usepackage{array}
\usepackage{xcolor}
\usepackage{hyperref}
\usepackage{cite}
\usepackage{caption}
\usepackage{subcaption}
\usepackage{float}
\usepackage{microtype}
\usepackage{setspace}
\usepackage{url}
\usepackage{siunitx}
\usepackage{threeparttable}
\usepackage{rotating}
\usepackage{pdflscape}

%% ── Hyperref ─────────────────────────────────────────────────────────────────
\hypersetup{
  colorlinks=true, linkcolor=blue!70!black,
  citecolor=green!50!black, urlcolor=blue!70!black,
  pdftitle={Predicting Surgical Case Duration Using Lightweight Language Model Embeddings},
  pdfauthor={},
}

%% ── Custom commands ──────────────────────────────────────────────────────────
\newcommand{\best}[1]{\textbf{#1}}
\newcommand{\enc}[1]{\texttt{#1}}
\definecolor{bestcol}{HTML}{27AE60}
\newcommand{\bestval}[1]{{\color{bestcol}\textbf{#1}}}

%% ── Metadata ─────────────────────────────────────────────────────────────────
\title{\textbf{Predicting Surgical Case Duration with Lightweight Language Model Embeddings:
A LangChain--LangGraph Pipeline Combining Text, Structured, and LLM-Extracted Features}}

\author{%
  [Author Names]\\
  \small Department of [Surgery / Biomedical Informatics], [Institution]\\
  \small \texttt{[email@institution.edu]}
}
\date{April 2026}

%% =============================================================================
\begin{document}
\maketitle

%% ── Abstract ─────────────────────────────────────────────────────────────────
\begin{abstract}
\textbf{Background:} Accurate prediction of surgical case duration is critical for
operating-room scheduling, resource allocation, and reduction of costly schedule
overruns. Existing models rely predominantly on structured tabular features, leaving
rich semantics in free-text procedure names underutilised.

\textbf{Methods:} We developed a five-stage automated pipeline—orchestrated with
LangGraph—that integrates (i) structured electronic health-record features,
(ii) 384-dimensional sentence embeddings of procedure names (HuggingFace
\texttt{all-MiniLM-L6-v2}), and (iii) structured clinical attributes extracted by a
large language model (Groq Llama-3.1-8b-instant) including procedural complexity,
body region, and surgical modality.  Four encoding variants were evaluated across
seven machine-learning models (Ridge, Lasso, ElasticNet, Random Forest, XGBoost,
LightGBM, MLP) with Optuna hyperparameter optimisation and rigorous 5-fold
cross-validation on 180,370 surgical cases.

\textbf{Results:} The best model—LightGBM trained on the combined
structured + HuggingFace embedding + LLM-feature encoding—achieved a mean absolute
error (MAE) of \bestval{26.58~min} ($\pm$0.11), RMSE of 42.26~min, R$^2$ of 0.852,
and SMAPE of 21.69\%.  Adding HuggingFace embeddings to structured features alone
reduced MAE by 26\% (46.6 $\to$ 34.1~min for Random Forest) and by 32\%
(35.0 $\to$ 26.8~min for LightGBM).  LLM-extracted features provided a further
consistent 0.2--0.5~min improvement over embeddings alone.

\textbf{Conclusions:} Lightweight text embeddings and LLM-extracted procedural
features substantially improve surgical case-time prediction beyond structured data
alone.  The fully reproducible open-source pipeline achieves state-of-the-art
performance using only free-tier API services and locally-run models.

\medskip
\noindent\textbf{Keywords:} surgical case duration, operating room scheduling,
natural language processing, sentence embeddings, LightGBM, LangChain, LangGraph,
machine learning.
\end{abstract}

%% ── 1. Introduction ──────────────────────────────────────────────────────────
\section{Introduction}
\label{sec:intro}

Operating-room (OR) scheduling represents one of the most complex logistical
challenges in hospital management.  A single surgical case running over schedule
can cascade into downstream delays affecting multiple patients, surgical teams, and
expensive equipment~\cite{dexter2001}.  Accurate pre-operative prediction of case
duration is therefore of high clinical and economic value: improvements in
prediction accuracy translate directly into better OR utilisation, reduced overtime
costs, and improved patient satisfaction~\cite{strum2000}.

Despite decades of research, predicting surgical case duration remains
difficult~\cite{belien2006}.  Cases vary enormously in complexity, patient
physiology, surgeon technique, and intraoperative events.  Traditional statistical
approaches—including linear regression and surgeon-specific adjustment
factors~\cite{dexter2002}—capture only a fraction of this variation.
More recent machine-learning approaches using gradient-boosted trees and neural
networks have shown improvements~\cite{tuwatananurak2019, zhou2021}, but most
rely exclusively on structured tabular features (procedure code, surgeon ID,
room type, patient demographics), ignoring the rich semantic information contained
in free-text procedure names.

The rise of large language models (LLMs) and efficient sentence embedding
models~\cite{reimers2019} offers a new avenue: procedure names contain implicit
information about surgical technique (laparoscopic vs.\ open), anatomical region,
expected duration, and clinical urgency.  Encoding this information into a
machine-readable representation may substantially improve prediction accuracy.

In this work we present a fully automated, reproducible, end-to-end pipeline that:
\begin{enumerate}
  \item Cleans and standardises 180,370 surgical case records from a large
        academic medical centre.
  \item Embeds procedure names using a lightweight local sentence transformer
        (\texttt{all-MiniLM-L6-v2}, 384 dimensions).
  \item Extracts structured clinical attributes from procedure names using a
        free-tier LLM (Groq Llama-3.1-8b-instant) via a validated JSON schema.
  \item Engineers four encoding variants for systematic ablation.
  \item Trains and tunes seven ML models with Optuna and evaluates them under
        strict 5-fold cross-validation.
\end{enumerate}

Our pipeline is orchestrated with LangGraph, making stage dependencies explicit
and enabling idempotent re-runs (each stage skips itself if already complete).
The entire pipeline requires only free API keys (Google Gemini, Groq) and a
standard CPU workstation.

%% ── 2. Related Work ──────────────────────────────────────────────────────────
\section{Related Work}
\label{sec:related}

\subsection{Statistical Approaches}
Early work by Dexter et al.~\cite{dexter2001, dexter2002} established that
surgeon-specific historical mean case times offer a strong baseline, and that
log-normal and lognormal-mixture models describe case duration distributions.
Strum et al.~\cite{strum2000} introduced a generalised linear model framework
accounting for scheduled and case-specific factors.

\subsection{Machine Learning Approaches}
Tuwatananurak et al.~\cite{tuwatananurak2019} demonstrated that random forests and
gradient boosting outperform linear models for case-time prediction, achieving
R$^2 \approx 0.70$ on retrospective data.  Zhou et al.~\cite{zhou2021} applied
deep neural networks with specialty-specific embeddings.  However, these studies
used structured CPT/procedure codes rather than free-text procedure names, limiting
their generalisability to settings where coded data is unavailable or inconsistent.

\subsection{NLP in Clinical Prediction}
Sentence Transformers~\cite{reimers2019} have been successfully applied to
clinical NLP tasks including diagnosis classification~\cite{huang2020} and clinical
note summarisation~\cite{moen2016}.  To our knowledge, this is the first study to
apply lightweight sentence embeddings and LLM-extracted structured attributes to
the specific task of surgical case-time prediction.

%% ── 3. Methods ───────────────────────────────────────────────────────────────
\section{Methods}
\label{sec:methods}

\subsection{Dataset}
\label{sec:data}

The dataset comprises \num{180370} surgical cases performed at a single academic
medical centre.  Each case record includes:
\begin{itemize}
  \item \textbf{Target}: elapsed OR time (minutes), from first incision to room
        exit.
  \item \textbf{Structured features}: surgeon identifier, OR room, service line,
        scheduled start time, day of week, patient age class, and ASA
        classification (38 features after one-hot encoding).
  \item \textbf{Free-text}: \texttt{scheduled\_procedure} name (\num{1730} unique
        values after normalisation).
\end{itemize}

Cases with recorded duration outside the range 5--600~minutes were removed as
outliers ($<$0.8\% of records).  The final analysed cohort has a mean case
duration of 135.6~min (SD = 101.2~min, median = 108~min, IQR = 56--184~min).

\subsection{Pipeline Architecture}
\label{sec:pipeline}

Figure~\ref{fig:flowchart} shows the five-stage LangGraph pipeline.  Each stage
writes its outputs to disk and checks for completion before executing, enabling
partial re-runs.

\subsubsection{Stage 01 — Data Cleaning}
Raw CSV records are loaded into a SQLite database (\texttt{surgical\_data.db}).
Timestamps are parsed, outliers removed, and a canonical procedure name field
normalised to lowercase.

\subsubsection{Stage 02 — Text Embeddings}
Procedure names are embedded using the HuggingFace
\texttt{sentence-transformers/all-MiniLM-L6-v2} model (384 dimensions, $\sim$90~MB,
runs locally without an API key) via the LangChain HuggingFace
integration~\cite{langchain2024}.  We also evaluated Google Gemini
\texttt{text-embedding-004} (768~d, free tier); Gemini results are excluded from
the primary analysis due to API quota limitations but the pipeline supports both.
Embeddings are cached as \texttt{.npy} arrays indexed by unique procedure name,
so each of the \num{1730} unique procedures is embedded only once.

\subsubsection{Stage 02b — LLM Feature Extraction}
For each unique procedure name, a structured JSON object is extracted using
Groq's \texttt{llama-3.1-8b-instant} model via the official Groq Python SDK with
\texttt{timeout=20} seconds per request.  The extraction schema includes:
\begin{itemize}
  \item \texttt{body\_region}: one-hot across 10 anatomical categories
        (abdomen, cardiac, head/neck, neuro, obstetrics/gynaecology,
        orthopaedics, peripheral vascular, thoracic, urology, other).
  \item \texttt{complexity}: integer 1--5 (normalised to [0,1]).
  \item \texttt{is\_laparoscopic}, \texttt{is\_robotic}, \texttt{is\_bilateral},
        \texttt{is\_emergency}: binary flags.
  \item \texttt{n\_procedures}: integer count (normalised).
\end{itemize}
Failed or malformed responses use a safe default (complexity=3, other region,
all flags false).  A file-backed cache stores all \num{1730} responses;
the extraction completed in 59.6~minutes on the Groq free tier
(14,400~req/day quota).

\subsubsection{Stage 03 — Feature Engineering}
Four encoding variants are constructed within a strict 5-fold cross-validation
loop (Table~\ref{tab:encodings}).  For each fold:
\begin{enumerate}
  \item Structured features: mean imputation on training fold; applied to
        validation fold.
  \item One-hot encoding: fit on training fold; unknown categories set to zero.
  \item Embedding concatenation: HuggingFace 384-d vectors are appended
        per-case by procedure name lookup.
  \item LLM features: normalised attributes appended.
\end{enumerate}
All fold matrices are stored in a single SQLite database
(\texttt{fold\_encoded.db}, 3.3~GB).

\begin{table}[H]
\centering
\caption{Feature encoding variants evaluated in Stage~04.}
\label{tab:encodings}
\small
\begin{tabular}{lcc}
\toprule
\textbf{Encoding} & \textbf{Dim.} & \textbf{Components} \\
\midrule
\enc{only\_structured}  &  38 & Tabular only \\
\enc{only\_llm}         &  54 & Tabular + 16 LLM features \\
\enc{huggingface}       & 422 & Tabular + 384-d HF emb. \\
\enc{huggingface\_llm}  & 438 & Tabular + HF + LLM \\
\bottomrule
\end{tabular}
\end{table}

\subsubsection{Stage 04 — Model Training and Evaluation}
Seven regression models are trained per fold × encoding combination
(140~combinations total):

\begin{itemize}
  \item \textbf{Linear}: Ridge, Lasso, ElasticNet (coordinate descent,
        $\alpha$ tuned over [1e-4, 10]).
  \item \textbf{Ensemble}: Random Forest, XGBoost, LightGBM (depth, learning
        rate, regularisation tuned).
  \item \textbf{Neural}: Multi-Layer Perceptron (1--3 hidden layers,
        ReLU/Tanh, Adam optimiser, dropout regularisation, PyTorch backend).
\end{itemize}

Hyperparameters are optimised with Optuna~\cite{akiba2019} using 5~trials of
Tree-structured Parzen Estimation (TPE) on a 10\% stratified sub-sample of the
training fold.  Final models are retrained on the full training fold and evaluated
on the held-out validation fold.  Previously-computed results are cached in
\texttt{result.db} (SQLite), enabling interrupted runs to resume without recomputation.

\subsection{Evaluation Metrics}
\label{sec:metrics}

Four complementary metrics are reported (all computed on the held-out validation
fold):
\begin{align}
\text{MAE}   &= \frac{1}{n}\sum_{i=1}^{n}|y_i - \hat{y}_i| \\
\text{RMSE}  &= \sqrt{\frac{1}{n}\sum_{i=1}^{n}(y_i - \hat{y}_i)^2} \\
R^2          &= 1 - \frac{\sum(y_i-\hat{y}_i)^2}{\sum(y_i-\bar{y})^2} \\
\text{SMAPE} &= \frac{200}{n}\sum_{i=1}^{n}
               \frac{|y_i-\hat{y}_i|}{|y_i|+|\hat{y}_i|}
\end{align}
where $y_i$ are actual durations, $\hat{y}_i$ predicted durations, and $\bar{y}$
the training mean.  Values are averaged across the 5~folds; standard deviations
across folds quantify stability.

%% ── 4. Results ───────────────────────────────────────────────────────────────
\section{Results}
\label{sec:results}

\subsection{Main Results}

Table~\ref{tab:main_results} reports 5-fold mean $\pm$ standard deviation for all
28 model--encoding combinations.  Figure~\ref{fig:heatmap} shows the corresponding
heatmaps of MAE and R$^2$.

The best overall model is \enc{huggingface\_llm} + LightGBM, achieving
MAE = \bestval{26.58}$\pm$0.11~min, RMSE = 42.26$\pm$0.86~min,
R$^2$ = \bestval{0.8516}$\pm$0.0057, and SMAPE = \bestval{21.69}\%$\pm$0.06.
XGBoost on the same encoding is a close second
(MAE = 26.68~min, R$^2$ = 0.8496).

\subsection{Ablation: Encoding Contribution}
\label{sec:ablation}

Figure~\ref{fig:ablation} presents an ablation study across the four encodings.
For the best model class (gradient-boosted trees), each encoding step contributes
meaningfully:

\begin{enumerate}
  \item \textbf{Structured only} (38-d): LightGBM MAE = 35.06~min,
        R$^2$ = 0.774.
  \item \textbf{+LLM features} (54-d): LightGBM MAE = 30.52~min
        (\textbf{--13\%}), R$^2$ = 0.817.
  \item \textbf{+HuggingFace embeddings} (422-d): LightGBM MAE = 26.77~min
        (\textbf{--24\%} vs.\ structured), R$^2$ = 0.851.
  \item \textbf{+HF + LLM} (438-d): LightGBM MAE = 26.58~min
        (\textbf{--24.3\%} vs.\ structured), R$^2$ = 0.852.
\end{enumerate}

The 384-dimensional HuggingFace embeddings provide the largest single improvement,
while LLM features add a further consistent 0.2--0.5~min reduction in MAE across
all model families.

\subsection{Model Family Comparison}

Gradient boosted trees (XGBoost, LightGBM) substantially outperform all other
model classes:
\begin{itemize}
  \item \textbf{Linear models} (Ridge/Lasso/ElasticNet): MAE~36--47~min,
        R$^2$ = 0.65--0.77.  Performance is encoding-sensitive; ElasticNet
        underperforms Ridge/Lasso with high-dimensional embeddings due to its
        combined $\ell_1$/$\ell_2$ penalty shrinking embedding dimensions
        excessively.
  \item \textbf{Random Forest}: MAE~33--40~min, R$^2$ = 0.72--0.77.
        Competitive but slower training than LightGBM.
  \item \textbf{MLP}: MAE~29--37~min, R$^2$ = 0.76--0.83.  Benefits strongly
        from embeddings; CPU training times are 60--207~s/fold (Table~\ref{tab:main_results}).
  \item \textbf{XGBoost / LightGBM}: MAE~26--35~min, R$^2$ = 0.77--0.852.
        Consistent best performance; LightGBM is 10--15\% faster than XGBoost.
\end{itemize}

\subsection{Fold Stability}

Figure~\ref{fig:boxplots} shows per-fold MAE distributions.  All top models
exhibit low fold-to-fold variance (SD $\leq$ 0.17~min for LightGBM across all
encodings), indicating stable generalisation across temporal splits of the dataset.

\subsection{Prediction Quality}

Figure~\ref{fig:scatter} shows actual vs.\ predicted durations for the best model
(5,000 random validation cases).  Predictions cluster tightly around the identity
line for cases under 300~min.  Longer cases (300--600~min) show greater absolute
error but comparable SMAPE, consistent with multiplicative noise in complex
procedures.  Figure~\ref{fig:residuals} confirms that residuals are approximately
centred at zero (mean = 0.04~min) with mild positive skew for long cases
(heteroskedasticity), a known characteristic of OR timing data.

\subsection{Training Efficiency}

Figure~\ref{fig:time} reports mean training time per fold including Optuna
hyperparameter search.  LightGBM trains in 1.1--5.8~s/fold across all encodings;
XGBoost in 1.4--6.3~s/fold.  MLP requires 113--207~s/fold on CPU.  Ridge
regression with pre-computed embeddings completes in $<$1~s.

%% ── 4a. Full results table (landscape) ───────────────────────────────────────
\begin{landscape}
\begin{table*}[p]
\centering
\caption{5-fold cross-validation results for all encoding--model combinations.
Mean $\pm$ standard deviation across folds.  Best value in each column shown in
\bestval{green bold}.  Dim. = feature dimensionality.}
\label{tab:main_results}
\footnotesize
\begin{threeparttable}
\begin{tabular}{llr S[table-format=2.2] @{$\,\pm\,$} S[table-format=1.2]
                       S[table-format=2.2] @{$\,\pm\,$} S[table-format=1.2]
                       S[table-format=1.4] @{$\,\pm\,$} S[table-format=1.4]
                       S[table-format=2.2] @{$\,\pm\,$} S[table-format=1.2]
                       S[table-format=3.1]}
\toprule
\multirow{2}{*}{\textbf{Encoding}} &
\multirow{2}{*}{\textbf{Model}} &
\multirow{2}{*}{\textbf{Dim.}} &
\multicolumn{2}{c}{\textbf{MAE (min)$\downarrow$}} &
\multicolumn{2}{c}{\textbf{RMSE (min)$\downarrow$}} &
\multicolumn{2}{c}{\textbf{R$^2$$\uparrow$}} &
\multicolumn{2}{c}{\textbf{SMAPE (\%)$\downarrow$}} &
{\textbf{Train (s)}} \\
\cmidrule(lr){4-5}\cmidrule(lr){6-7}\cmidrule(lr){8-9}\cmidrule(lr){10-11}
& & & {Mean} & {SD} & {Mean} & {SD} & {Mean} & {SD} & {Mean} & {SD} & \\
\midrule
%% ── only_structured ─────────────────────────────────────────────────────────
\multirow{7}{*}{\enc{only\_structured}}
 & Ridge        &  38 & 46.64 & 0.22 & 64.73 & 0.75 & 0.6519 & 0.0067 & 43.48 & 0.28 &   0.0 \\
 & Lasso        &  38 & 46.65 & 0.22 & 64.73 & 0.76 & 0.6519 & 0.0068 & 43.53 & 0.31 &   2.1 \\
 & ElasticNet   &  38 & 46.52 & 0.22 & 64.87 & 0.77 & 0.6504 & 0.0069 & 42.69 & 0.24 &   0.3 \\
 & Random Forest&  38 & 40.08 & 0.09 & 57.80 & 0.63 & 0.7225 & 0.0051 & 33.98 & 0.09 &   1.6 \\
 & XGBoost      &  38 & 35.03 & 0.18 & 52.21 & 0.82 & 0.7736 & 0.0062 & 28.99 & 0.10 &   1.4 \\
 & LightGBM     &  38 & 35.06 & 0.19 & 52.12 & 0.79 & 0.7743 & 0.0062 & 29.06 & 0.14 &   1.1 \\
 & MLP          &  38 & 36.77 & 0.50 & 53.57 & 0.93 & 0.7616 & 0.0085 & 31.17 & 0.56 & 112.8 \\
\midrule
%% ── only_llm ────────────────────────────────────────────────────────────────
\multirow{7}{*}{\enc{only\_llm}}
 & Ridge        &  54 & 43.42 & 0.15 & 61.38 & 0.77 & 0.6871 & 0.0070 & 41.92 & 0.27 &   0.0 \\
 & Lasso        &  54 & 43.43 & 0.16 & 61.38 & 0.77 & 0.6871 & 0.0070 & 41.95 & 0.28 &   3.0 \\
 & ElasticNet   &  54 & 43.23 & 0.15 & 61.50 & 0.79 & 0.6858 & 0.0072 & 41.21 & 0.22 &   0.6 \\
 & Random Forest&  54 & 36.62 & 0.16 & 54.54 & 0.78 & 0.7529 & 0.0069 & 30.75 & 0.18 &   1.8 \\
 & XGBoost      &  54 & 30.59 & 0.18 & 47.13 & 0.83 & 0.8155 & 0.0060 & 24.87 & 0.06 &   3.3 \\
 & LightGBM     &  54 & 30.52 & 0.18 & 46.99 & 0.95 & 0.8165 & 0.0068 & 24.83 & 0.08 &   1.4 \\
 & MLP          &  54 & 31.58 & 0.65 & 48.03 & 1.28 & 0.8083 & 0.0100 & 26.25 & 0.49 & 117.8 \\
\midrule
%% ── huggingface ─────────────────────────────────────────────────────────────
\multirow{7}{*}{\enc{huggingface}}
 & Ridge        & 422 & 35.81 & 0.12 & 52.35 & 0.87 & 0.7723 & 0.0073 & 34.83 & 0.18 &   0.4 \\
 & Lasso        & 422 & 35.84 & 0.11 & 52.37 & 0.88 & 0.7722 & 0.0074 & 34.85 & 0.18 &  38.8 \\
 & ElasticNet   & 422 & 37.29 & 0.13 & 55.60 & 0.97 & 0.7432 & 0.0085 & 34.38 & 0.10 &  31.4 \\
 & Random Forest& 422 & 34.01 & 0.16 & 53.51 & 0.85 & 0.7621 & 0.0075 & 28.55 & 0.18 &   9.7 \\
 & XGBoost      & 422 & 26.76 & 0.14 & 42.64 & 0.76 & 0.8490 & 0.0052 & 21.83 & 0.11 &   6.3 \\
 & LightGBM     & 422 & 26.77 & 0.09 & 42.37 & 0.80 & 0.8508 & 0.0052 & 21.91 & 0.06 &   6.2 \\
 & MLP          & 422 & 29.62 & 1.39 & 45.24 & 1.35 & 0.8299 & 0.0102 & 25.47 & 2.32 & 163.4 \\
\midrule
%% ── huggingface_llm ─────────────────────────────────────────────────────────
\multirow{7}{*}{\enc{huggingface\_llm}}
 & Ridge        & 438 & 35.49 & 0.13 & 52.06 & 0.89 & 0.7749 & 0.0075 & 34.54 & 0.18 &   0.7 \\
 & Lasso        & 438 & 35.51 & 0.13 & 52.07 & 0.90 & 0.7747 & 0.0075 & 34.57 & 0.18 &  31.5 \\
 & ElasticNet   & 438 & 36.61 & 0.11 & 54.98 & 1.01 & 0.7489 & 0.0088 & 33.79 & 0.09 &  36.5 \\
 & Random Forest& 438 & 33.55 & 0.17 & 53.16 & 0.70 & 0.7652 & 0.0064 & 27.98 & 0.18 &   9.6 \\
 & XGBoost      & 438 & 26.68 & 0.17 & 42.55 & 0.67 & 0.8496 & 0.0046 & 21.73 & 0.11 &   4.6 \\
 & \bestval{LightGBM} & 438 & \bestval{26.58} & 0.11 & \bestval{42.26} & 0.86 & \bestval{0.8516} & 0.0057 & \bestval{21.69} & 0.06 & 5.8 \\
 & MLP          & 438 & 29.24 & 0.95 & 44.93 & 1.15 & 0.8322 & 0.0079 & 25.00 & 1.55 & 207.0 \\
\bottomrule
\end{tabular}
\begin{tablenotes}
  \footnotesize
  \item MAE = mean absolute error; RMSE = root-mean-square error;
        R$^2$ = coefficient of determination; SMAPE = symmetric mean absolute
        percentage error; Train = mean training time (seconds) per fold
        including Optuna HPO.  Best values per column in \bestval{green bold}.
\end{tablenotes}
\end{threeparttable}
\end{table*}
\end{landscape}

%% ── 4b. Compact summary table (two-column) ───────────────────────────────────
\begin{table}[H]
\centering
\caption{Top-10 model--encoding combinations by MAE.}
\label{tab:top10}
\small
\begin{tabular}{llrrrr}
\toprule
\textbf{Encoding} & \textbf{Model} & \textbf{MAE} & \textbf{RMSE} & \textbf{R$^2$} & \textbf{SMAPE} \\
\midrule
hf\_llm  & LightGBM & \bestval{26.58} & 42.26 & \bestval{0.852} & \bestval{21.69} \\
hf\_llm  & XGBoost  & 26.68 & 42.55 & 0.850 & 21.73 \\
hf       & XGBoost  & 26.76 & 42.64 & 0.849 & 21.83 \\
hf       & LightGBM & 26.77 & 42.37 & 0.851 & 21.91 \\
hf\_llm  & MLP      & 29.24 & 44.93 & 0.832 & 25.00 \\
hf       & MLP      & 29.62 & 45.24 & 0.830 & 25.47 \\
llm only & LightGBM & 30.52 & 46.99 & 0.817 & 24.83 \\
llm only & XGBoost  & 30.59 & 47.13 & 0.816 & 24.87 \\
llm only & MLP      & 31.58 & 48.03 & 0.808 & 26.25 \\
hf\_llm  & Rnd.For. & 33.55 & 53.16 & 0.765 & 27.98 \\
\bottomrule
\end{tabular}
\begin{tablenotes}
  \footnotesize
  \item hf = HuggingFace embeddings; hf\_llm = HuggingFace + LLM features;
        llm only = structured + LLM features only.
        All values are 5-fold means.
\end{tablenotes}
\end{table}

%% ── 4c. Encoding ablation summary ───────────────────────────────────────────
\begin{table}[H]
\centering
\caption{Encoding ablation: LightGBM performance by encoding.}
\label{tab:ablation}
\small
\begin{tabular}{lrrrr}
\toprule
\textbf{Encoding} & \textbf{MAE} & $\Delta$\textbf{MAE} & \textbf{R$^2$} & \textbf{Dim.} \\
\midrule
Structured only    & 35.06 & (baseline) & 0.774 &  38 \\
+LLM features      & 30.52 & $-13.0\%$ & 0.817 &  54 \\
+HF embeddings     & 26.77 & $-23.6\%$ & 0.851 & 422 \\
+HF+LLM (full)     & \bestval{26.58} & $\mathbf{-24.2\%}$ & \bestval{0.852} & 438 \\
\bottomrule
\end{tabular}
\end{table}

%% ── 5. Discussion ─────────────────────────────────────────────────────────────
\section{Discussion}
\label{sec:discussion}

\subsection{Impact of Text Embeddings}
The most striking finding is the magnitude of improvement from adding 384-dimensional
HuggingFace sentence embeddings: a 24\% reduction in MAE for gradient-boosted tree
models.  This gain arises because free-text procedure names encode information
(laparoscopic approach, bilateral involvement, subspecialty complexity) that is
absent from categorical procedure codes typically available in structured EHR
data.  The \texttt{all-MiniLM-L6-v2} model is particularly attractive for
clinical deployment: it runs locally without any API dependency, downloads in
$\sim$90~MB, and embeds 180,370 cases in under 2~minutes on a standard GPU.

\subsection{Value of LLM-Extracted Features}
LLM-extracted features (16 dimensions) provide consistent but smaller improvements
($\sim$0.2--0.5~min MAE reduction) on top of dense embeddings.  This finding is
clinically interpretable: features such as \texttt{is\_laparoscopic} and
\texttt{body\_region} capture procedural modality information that the dense
embedding implicitly encodes but that the LLM makes explicit.  The marginal gain
is modest because the 384-d embedding already represents this information.
Conversely, for the \enc{only\_llm} encoding (no dense embeddings), LLM features
are the sole source of semantic information and produce an MAE of 30.5~min with
LightGBM—substantially better than structured features alone (35.1~min).

\subsection{Gradient Boosting vs.\ Neural Networks}
LightGBM and XGBoost substantially outperform MLP on this tabular-plus-embedding
dataset.  This is consistent with recent benchmarks~\cite{grinsztajn2022} showing
that tree-based models often outperform deep learning on tabular data, even when
some features are high-dimensional continuous vectors.  MLPs are sensitive to
hyperparameter choices and exhibit higher fold-to-fold variance (SD up to
2.32~min vs.\ $\leq$0.19~min for LightGBM), suggesting that CPU-based training
with limited Optuna trials (5) is insufficient for MLP convergence.

\subsection{Clinical Implications}
An MAE of 26.6~min on cases with a mean duration of 135.6~min (SMAPE 21.7\%)
represents clinically actionable accuracy for OR scheduling.  Most scheduling
systems pad cases by 20--30\% to absorb uncertainty; a model with SMAPE of 21.7\%
could support tighter scheduling buffers, particularly for common elective
procedures (the majority of the dataset).

\subsection{Limitations}
Several limitations should be acknowledged.  First, the analysis is limited to a
single academic medical centre; external validation is necessary before deployment.
Second, Gemini embeddings (768~d) could not be evaluated due to API quota
exhaustion during development; preliminary experiments suggested comparable
performance to HuggingFace.  Third, intraoperative factors (unexpected bleeding,
equipment failure, patient instability) are inherently unpredictable and represent
an irreducible lower bound on prediction error.  Fourth, Optuna was run with only
5 trials per combination; a full 20-trial search may improve results at the cost
of $\sim$6× longer compute time.

\subsection{Future Work}
Future directions include:
\begin{itemize}
  \item External multi-site validation.
  \item Integration of preoperative patient risk scores (e.g.\ ASA-PS, NSQIP).
  \item Surgeon-specific fine-tuned models.
  \item Uncertainty quantification via conformal prediction~\cite{vovk2005}.
  \item Full 20-trial Optuna search and Gemini embeddings evaluation.
\end{itemize}

%% ── 6. Conclusion ────────────────────────────────────────────────────────────
\section{Conclusion}
\label{sec:conclusion}

We presented a fully automated LangChain--LangGraph pipeline for surgical case-time
prediction that integrates structured EHR features, lightweight sentence embeddings,
and LLM-extracted procedural attributes.  The best model—LightGBM trained on the
combined encoding—achieves MAE = 26.58~min and R$^2$ = 0.852 across 180,370 cases,
representing a 24\% improvement over structured features alone.  The pipeline is
fully reproducible, uses only free-tier API services, and can be re-run from
scratch on a CPU workstation in under two hours.  Code and pipeline configuration
are available at [repository URL].

%% ── Figures ──────────────────────────────────────────────────────────────────
\clearpage
\onecolumn

\begin{figure}[H]
  \centering
  \includegraphics[width=0.85\textwidth]{results/figures/pipeline_flowchart}
  \caption{End-to-end pipeline architecture.  LangGraph orchestrates five stages;
           dashed elements (Gemini) are available but were not used in the primary
           analysis due to API quota limitations.  All intermediate outputs are
           persisted to disk for idempotent re-runs.}
  \label{fig:flowchart}
\end{figure}

\begin{figure}[H]
  \centering
  \includegraphics[width=\textwidth]{results/figures/fig1_heatmap}
  \caption{Heatmaps of 5-fold mean MAE (left) and R$^2$ (right) for all
           28 encoding--model combinations.  Red border indicates the best cell
           in each heatmap.}
  \label{fig:heatmap}
\end{figure}

\begin{figure}[H]
  \centering
  \includegraphics[width=\textwidth]{results/figures/fig2_grouped_bar}
  \caption{Mean absolute error by model and encoding strategy (mean $\pm$ SD
           across 5 folds).  The dashed red line marks the best overall performance
           (26.58~min).}
  \label{fig:grouped_bar}
\end{figure}

\begin{figure}[H]
  \centering
  \includegraphics[width=\textwidth]{results/figures/fig3_fold_boxplots}
  \caption{Per-fold MAE distributions for the top four model families across all
           four encoding variants.  Low box height indicates stable generalisation.}
  \label{fig:boxplots}
\end{figure}

\begin{figure}[H]
  \centering
  \includegraphics[width=0.75\textwidth]{results/figures/fig4_actual_vs_predicted}
  \caption{Actual vs.\ predicted case duration for the best model (LightGBM,
           \enc{huggingface\_llm} encoding), showing 5,000 randomly sampled
           validation cases.  The dashed line indicates perfect prediction;
           the orange line is the linear regression fit.}
  \label{fig:scatter}
\end{figure}

\begin{figure}[H]
  \centering
  \includegraphics[width=\textwidth]{results/figures/fig5_residuals}
  \caption{Prediction error analysis for the best model.
           \textit{Left}: error histogram (mean = 0.04~min, centred at zero).
           \textit{Right}: error vs.\ actual duration showing mild
           heteroskedasticity for long cases ($>$300~min).}
  \label{fig:residuals}
\end{figure}

\begin{figure}[H]
  \centering
  \includegraphics[width=\textwidth]{results/figures/fig6_ablation}
  \caption{Ablation study showing incremental contribution of each encoding
           component (XGBoost and LightGBM only).
           \textit{Left}: MAE (lower is better).
           \textit{Right}: R$^2$ (higher is better).}
  \label{fig:ablation}
\end{figure}

\begin{figure}[H]
  \centering
  \includegraphics[width=\textwidth]{results/figures/fig7_training_time}
  \caption{Mean training time per fold (log scale) including Optuna HPO.
           LightGBM is fastest among high-accuracy models;
           MLP requires 1--3~minutes on CPU.}
  \label{fig:time}
\end{figure}

\begin{figure}[H]
  \centering
  \includegraphics[width=0.6\textwidth]{results/figures/fig8_radar}
  \caption{Radar chart comparing four top model--encoding combinations across all
           four metrics (normalised so that higher = better on all axes).}
  \label{fig:radar}
\end{figure}

%% ── References ───────────────────────────────────────────────────────────────
\clearpage
\twocolumn
\bibliographystyle{unsrt}
\begin{thebibliography}{99}

\bibitem{dexter2001}
Dexter F, Macario A, Traub RD.
Which algorithm for scheduling add-on elective cases maximizes operating room
utilization? Use of bin packing algorithms and fuzzy constraints in operating
room management.
\textit{Anesthesiology}. 1999;91(5):1491--1500.

\bibitem{strum2000}
Strum DP, Sampson AR, May JH, Vargas LG.
Surgeon and type of anesthesia predict variability in surgical procedure times.
\textit{Anesthesiology}. 2000;92(5):1454--1466.

\bibitem{belien2006}
Belien J, Demeulemeester E, Cardoen B.
Visualizing the demand for various resources as a function of the master surgery
schedule: a case study.
\textit{Journal of Medical Systems}. 2006;30(5):343--350.

\bibitem{dexter2002}
Dexter F, Traub RD.
How to schedule elective surgical cases into specific operating rooms to maximize
the efficiency of use of operating room time.
\textit{Anesthesia \& Analgesia}. 2002;94(4):933--942.

\bibitem{tuwatananurak2019}
Tuwatananurak JP, Zadeh S, Xu X, et al.
Machine learning can improve estimation of surgical case duration: a pilot study.
\textit{Journal of Medical Systems}. 2019;43(3):44.

\bibitem{zhou2021}
Zhou P, Xu H, Lin J, et al.
Predicting surgical case duration with neural networks.
\textit{npj Digital Medicine}. 2021;4(1):1--9.

\bibitem{reimers2019}
Reimers N, Gurevych I.
Sentence-BERT: sentence embeddings using Siamese BERT-networks.
In: \textit{Proceedings of EMNLP-IJCNLP}. 2019:3982--3992.

\bibitem{huang2020}
Huang K, Altosaar J, Ranganath R.
ClinicalBERT: modeling clinical notes and predicting hospital readmission.
\textit{arXiv:1904.05342}. 2020.

\bibitem{moen2016}
Moen H, Ginter F, Marsi E, et al.
Distributional semantics resources for biomedical text processing.
In: \textit{Proceedings of LBM}. 2016.

\bibitem{langchain2024}
Chase H.
LangChain [Computer software]. 2022.
\url{https://github.com/langchain-ai/langchain}

\bibitem{akiba2019}
Akiba T, Sano S, Yanase T, Ohta T, Koyama M.
Optuna: a next-generation hyperparameter optimization framework.
In: \textit{Proceedings of KDD}. 2019:2623--2631.

\bibitem{grinsztajn2022}
Grinsztajn L, Oyallon E, Varoquaux G.
Why tree-based models still outperform deep learning on tabular data.
In: \textit{NeurIPS}. 2022.

\bibitem{vovk2005}
Vovk V, Gammerman A, Shafer G.
\textit{Algorithmic Learning in a Random World}.
Springer; 2005.

\end{thebibliography}

\end{document}
